% Oppgaven kan skrives på norsk eller engelsk. Hvis den skrives på engelsk må i så fall overskriftene oppdateres.
% bruk oppsett her
% Merk at 25% av karakteren på oppgaven er struktur og form og vi anbefaler derfor alle om følge denne malen for å unngå unødvendig trekk i karakter.

% Husk å kjøre pdflatex 2 ganger etter bibtex:
% pdflatex filnavn
% bibtex filnavn
% pdflatex filnavn
% pdflatex filnavn

% MERK: det er nok flere pakker som må innstalleres her!
\documentclass[11pt]{article} 
\usepackage[utf8]{inputenc} 
\usepackage{geometry} 
\geometry{a4paper} 
\usepackage{graphicx} 
\usepackage{fancyhdr}
\usepackage[english]{babel} % 'norsk' må endres til 'english' hvis du skal skrive på Engelsk
% Inkluder eventuelle pakker her

% START:  DISSE MÅ ENDRES
\newcommand{\tittel}{TITTEL} % Tittel velger du 
\newcommand{\temanr}{26} % referer til det som er oppgitt på wiki-siden
\newcommand{\tematittel}{Tracking and mass surveillance in the era of Big Data} % referer til det som er oppgitt på wiki-siden 
\newcommand{\kandidatnr}{julihau, olejfri} % komma-separert liste med kandidatnummer for alle forfatter VIKTIG: bruk kandidatnummer og ikke studentnummer siden oppgaven skal være anonym.
\newcommand{\kandidatnrforside}{julihau \\ olejfri}  %  liste med kandidatnummer for alle forfatter - her skal hver forfatter separeres med \\
\newcommand{\antallord}{ORD} % antall ord må fylles inn (trenger ikke være 100% korrekt men skal gi en indikasjon; flere måter texcount, kompiler til PDF og tell, kopier i word og bruk wordcount, +++,  
% Merk: lengden skal være mellom 5000 – 10000 ord 
% Man må påregne trekk i karakter dersom oppgaven er nevneverdig kortere eller lengre enn dette, men antall ord innenfor gitt intervall vil ikke ha betydning på karakter. Her bedømmes det utfra kvalitet og ikke kvantitet. 
% SLUTT:  DISSE MÅ ENDRES


% START: DETTE SKAL IKKE ENDRES
\pagestyle{fancy}
\fancyhf{}
\fancyhead[RO,RE]{\kandidatnr}
\fancyhead[LO,LE]{IN2120 -- midtsemestereksamen 2022}
\fancyfoot[C]{\thepage}
\begin{document} \thispagestyle{empty}
\begin{center}
\normalsize{\textit{Universitetet i Oslo -- IN2120 Informasjonssikkerhet-- 2022 -- midtsemestereksamen}} \\
\vspace{1cm}
\includegraphics[width=0.15\textwidth]{uio_logo.png} \\
\vspace{1cm}
\huge{\tittel} \\
\Large{\temanr} \\
\Large{\tematittel} \\
\vspace{1cm}
\Large{Kandidatnummer:} \\
\large{\kandidatnrforside} \\
\vspace{1cm}
\Large{Antall ord: \antallord} \\
\end{center}
\tableofcontents
\newpage
% SLUTT: DETTE SKAL IKKE ENDRES


% HER KOMMER SELVE TEKSTEN I OPPGAVEN

\section{Abstract}
How have legislations, attitudes towards and the consequences of tracking and mass surveillance changed over the course of time, taking into consideration the recent global pandemic and how data-tracking was used as a countermeasure. \\ \\

The practice of mass surveillance was first put to use as a result of the happenings of 9/11 2001 (Sinha, 2013) with the reasoning that it could be used to help prevent future terrorist attacks and the like. It was presented as a necessary evil in order to assist authorities in ensuring the safety of the people, and as such, a new normal emerged. By using massive sets of data collected from its netizens, new systems were created to predict anything from traffic jams to premediated crime (Zuboff, 2015).  \\ \\

The people’s willingness to volunteer their data likely stems from the theory that in order to feel a sense of safety from these frightening prospects, a person would rather sacrifice their own privacy than risk becoming a victim themselves (Zuboff, 2019). This would however turn out to be the beginning of something much bigger, namely the understanding of how much you could actually discern about a person by analyzing their data, and how big companies could profit from doing exactly that (Vargo \& Hopp, 2020).  \\ \\

There have been a few notable cases of tracking, manipulation and surveillance is years following the start of this new paradigm. One example being the Cambridge Analytica scandal, in which it was brought to light that Facebook had handed over personally identifiable information of more than 87 million users to to Cambridge Analytica. It would later be discovered that this information was used to influence the American presidential election of 2016 (Rathi, 2019). What was made clear following scandals such as this was that this new era carried with it the risk of private companies being able to control the masses (Karas, 2002), even to the point of national leaders being elected by their will. \\ \\

During the Covid-19 pandemic, multiple forms of mobile applications with the express purpose of tracking the spread of the virus were introduced. Both government-approved applications and ones developed by notorious private companies such as Apple and Google (Kaashoek \& Santillana, 2020) were put to use on an international level (Munzert, et al., 2021). This begs the question; “Where does user privacy begin and end?” How can we be sure that the data we are giving away, with the intent of benefitting the public good, is used in a manner that is effective or at all useful?” Could this be the start of an era where mass data-collection and surveillance are an everyday phenomenon? What data will be standard procedure to collect from users of mobile applications, and how could this information be used in the future? It would seem that it all boils down to the question; “In what ways are we allowing both authorities and private companies to influence our lives by sharing our data?”. 


%Innledningen skal gi en kort introduksjon til tema og motivasjon til oppgaven, og beskrive hva oppgaven skal besvare. Denne vil typisk være rundt en side (men dette er ikke et krav vil variere fra tema til tema).
%Sitering/referanser: 25% av karakteren er god bruk av kilder og referanser. Vi anbefaler bruk av passende sitering i teksten (se under) og vil anbefale at dere bruker bibtex.
% Dere kan bruke hvilken stil som dere ønsker og her har vi brukt en enkel tall-referanse (Word-malen bruker Harvard-stilen) - det er helt opp til dere hvilken stil som brukes.

%Eksempel referanse: \cite{josang21} (denne setningen er for illustrasjon og må fjernes)%

\section{Hovedkapittel}

%Besvarelsen må ha ett eller flere hovedkapitler.  Hvert hovedkapittel kan ha så mange delkapitler som ønskelig (så fremt det gir mening strukturmessig) og skal være nummert som gjort i denne malenn. Fjern/legg til etter behov. Anbefaler klipp-og-lim eksisterende kapittel/delkapittel for å få riktig nummerte overskrifter.

Due to the increased popularity of personal profiles on social media sites and the like, we find ourselves in a period of mass personalization (Mulvenna, et al., 2000). Many websites and applications like Facebook, Twitter and Tumblr can potentially create a comprehensive list of personal data such as demographic, behavioral, psychographic and lifestyle-data simply by using its intended functions (Russel, 2013) (Arya, et al., 2019). In addition to this mined data, many applications such as WhatsApp, Messenger and TikTok include a statement in their policies, informing us that by using their applications, the user is consenting to their communication being put through textual analysis and searched for keywords. These keywords are then used to personalize what advertisements are shown to the user and other users who display similar textual patterns (Vargo \& Hopp, 2020). 

According to (Trang, 2017), advertising strategies have been changed drastically as a result of the development of social media. Due to their 1.39 million active monthly users, Facebook has become one of the most common online marketing tools with 92\% of social marketing companies having chosen to use the platform due to their comparatively low prices and massive user base.  

In addition to this traditional form of advertisement, Facebook teamed up with IBM on May 6th, 2015 with the purpose of creating better advertising on what was then, and continues to be, the world’s largest social network. By combining Facebook’s targeting technology and IBM’s services for marketers, they created a new form of advertisement system where users’ interactions on collaborating marketers’ websites resulted in advertisements of viewed products or services being generated and put on their Facebook feed. An example of this would be if a user visited company X’s website and looked at a certain product. This viewing data would then be sent by the company to Facebook, and an ad of that product featuring an affiliate link to company X’s website would then later appear on the user’s Facebook feed (Trang, 2017).  

While this practice has been shown to have very high success rates compared to the more traditional methods of advertising, it has brought with it the consequence that selling these massive amounts of data to third parties has become a big business in the 21st century (Palos-Sanchez, et al., 2019). 

\subsection{Manipulation and Recruitment}
We have already discussed why using or selling user data for advertising purposes may be one of the main reasons for websites and applications to collect it, but taking into account the questions we posed earlier, there is another aspect to consider. Namely user data being used by potentially malicious parties wishing to manipulate or recruit the user.\\ \\ 

As previously mentioned, one notorious example of users can be manipulated as a result of their data being shared is the Cambridge Analytica scandal. \\ \\

Cambridge Analytica was officially an advertising company with both a commercial and a political wing. The political wing was claimed to “combine the predictive data analytics, behavioral sciences, and innovative ad tech into one award winning approach”. The method in practice involved making unique advertisements for many different categories of users which were grouped based on their analyzed traits. These analyzed traits were the results of a regression model that in the words of Professor David Sumpter: “Takes the data we already know about a person, and uses them to predict something we don’t know about him or her” (Sumpter, 2018). The idea behind it all was that different types of people will react more strongly to different forms of advertisement, be it appealing to the viewers sentimentality or by inciting feelings of anger with the opposition. By exposing Facebook users to these tailored ads, they intended to yield better results compared to the more traditional blanket advertisement where the same ad is shown to everyone.  
\\ \\
The effectiveness of this approach has been both exaggerated and underplayed by both the media, CEO Alexander Nix, and the whistleblower and ex-employee Christopher Wylie, but the fact remains that Cambridge Analytica, to an unknown extent, helped republican politician Ted Cruz get the second highest number of votes in the republican primary election of 2016 (Rathi, 2019). 
\\ \\
Despite the arguments of both sides regarding how effective or ineffective these kinds of regression models are, users of platforms such as Facebook are in fact having their data used in what has been called both “an ethical grey area” and “an attempt to manipulate voters by latching onto their vulnerabilities” (McCausland \& Schnecter, 2018). \\ \\

 

In addition to this form of manipulation, we can also refer to the research done by (Alava, et al., 2017) which speaks of a relatively recent, but growing, phenomenon often referred to as “incitement to radicalization towards violent extremism” or more simply put “violent radicalization”. The situation we are faced with according to their research is that this boom of social media has created new and easier ways for radical groups to identify and approach users that they believe could become allies to their cause. Taking into account that these groups, given sufficient funds could purchase potentially massive amounts of user information to aid them in their recruitment, puts to question whether the premise that the gathering and availability of personal data is in the best interest of public safety. \\ \\

 

Taking what we have discussed about the subject into account, we can say that by using applications and services like Facebook, Twitter and Snapchat, the user is either knowingly or unknowingly offering their data to the company and whomever they choose to share that data with. This data may be used against the user in the form of targeted advertising, manipulation or potentially tracking and surveillance depending on what data the application chooses to collect. One could argue that a common denominator that allowed, and still allows, these happenings to take place is the culture of mass data collection which, as previously mentioned, has created a thriving market for the commerce and utilization of user data.  \\ \\

\subsection{The COVID-19 pandemic}
The question of whether or not the handling and management of user data is a violation of user privacy has been the subject of many studies in this era of Big Data (Buchanan, et al., 2006). Something that brought more attention to the issue was however the introduction of the COVID-19 related tracking of citizens. \\ \\

Since the beginning of the virus-outbreak in 2020, the pandemic caused many systematic changes in organizational structures and society as a whole. Surveillance and control in the form of collecting data about location and interactions for its citizens was advised worldwide, which caused fear for the thought of a future where pandemic countermeasures could systematically invade user privacy (Ribeiro-Navarrete, et al., 2021).  \\ \\

After the recognition of the second wave of the COVID-19-pandemic, increasingly more countries chose to implement large-scale technology-based tracking measures in order to monitor and prevent the spread of the virus across the population. Despite the governments’ efforts, these technology-based tracking applications ultimately experienced a global failure to convince the public due to these applications’ perceived risk of infringing individual liberties and were met with very low levels of trust from the population (Georgieva, et al., 2021).  \\ \\

The findings in the article by (Hargittau, et al., 2020) describe the public distrust towards contact-tracing apps as being the product of a decade of media-covered security breaches and other violations of users’ privacy expectations with one notable example being ”Edward Snowden’s revelations of mass surveillance”, which is thought to have single-handedly caused a significant decrease of trust in the federal government. When trust in an institution decreases, people are more likely to turn to a competitor, but with research showing a steady decline of trust in technology companies in light of happenings such as the Cambridge Analytica scandal (Selinger \& Hartzog, 2016), these non-governmental contract-tracking apps are distrusted as well. \\ \\

 

This situation of distrust towards both governmental and privately owned contact-tracing apps has lead to a dilemma in which the government can chose to use other means in order to track the contagion of the virus, but at the cost of using their power to commandeer location and interaction information from cellular service providers. \\ \\

(skriv om denne her https://www.tandfonline.com/doi/full/10.1080/0960085X.2020.1802358) 

\subsection{How powerful is really more data for large corporations?}
The Central Limit Theorem posits that methods that apply to normal distributions work also for any other distribution because sums and averages are normal for a sufficiently large sample-size, $n = 50$ being large enough (Devore, 2018, p. 298-299). That is quite a large room of application with only 50 independent random variables from a sample. This gives an opportunity to study confidence intervals of the mean; and create test-statistics, calculate a p-value and do hypothesis testing of any distribution.\\ \\

Other times a larger sample is necessary; Say, when creating confidence intervals of variance (Devore, 2018, p. 410) because critical values to establish interval boundaries are so high; Or when achieving a high power to reject the null hypothesis when it is false (Devore, 2018, p. 468). Then again, is it important to know that something is 99.99\% likely opposed to 99.9\% or that the population mean is $\mu = 100$ and not $\mu = 100.1$? \\ \\

Machine learning on the other hand is a branch under artificial intelligence that seeks to build models algorithmically based on the sample data while in statistics known models are tested and checked for how well they perform in relation to the sample. In order to check how well a model performs, a statistician uses mathematical methods to create a model him/her-self. This is not necessary in machine learning where the power of computers are leveraged. It is interesting to see how powerful this approach has become. AI now outperforms humans in question-answering task (Kumar) and drives cars. It has been able to vastly expand on data science from what statistics has offered. \\ \\

Machine learning applications are hungry for data (Lipeles). As long as one has a lot of data, the neural architecture doesn't seem to matter too greatly. The precision between state-of-the-art natural language processing networks isn't huge (Liu, p. 8) and longer training times and larger data-sets may contribute more to performance then architecture (Liu, p. 7)

\subsection{Governmental data collection, tracking and surveillance}
Paraphrasing the words of (Cate, 2008), The United States government has long sought data about individuals for a wide variety of important public purposes, but due to the then time and cost-heavy processes yielding mostly useless data, this collecting of data was protected under The Supreme Courts concept of “practical obscurity”.  \\ \\

The concept stems from the case “U.S. Department of Justice v. Reporters Committee for Freedom of Press” of 1989, and in practice means that while some information is technically accessible to the public, the time and effort needed to acquire this information makes it somewhat protected from mass access. The example of this used in the Supreme court case was that whilst the rap sheets and court records of a person were technically available to any person with the time to dig through a county courthouse’s paper records in search of them, the high cost and low chance of acquiring them meant that the information had little chance of being compiled by the public (Reporters Committee for Freedom of the Press, 2009).  \\ \\

While this concept was deemed adequate protection of privacy back in 1989, with new technology, the situation has changed. Despite court records and arrest reports being available only as physical copies in most U.S. counties, corporate data brokers are free to collect this information on individuals from all over the country which they can then sell to government agencies, insurance companies and other third parties belonging to the private sector. A notable example of such a company is ChoicePoint which offers aggregate reports of a person’s line of credit collected from multiple public and private databases (Personal Finance, 2019). These reports can in some cases determine whether a person is considered qualified for a job, professional license, or insurance (Testa-Avila, 2008). \\ \\

While the government does not acquire this data directly, the Supreme Court has refused to restrict the government’s access to data held by third parties which they could do by extending the Fourth Amendment that was originally meant to ensure the right to privacy (Cate, 2008). As it currently stands, companies have common authority over their business records and can therefore consent to a government search or acquisition of their databases even if their users oppose it (Kerr, 2021). \\ \\

This enforces that by using applications, websites or services, the user allows both private sector parties and authorities access to their personal data which could be used for both surveillance and tracking. 

\subsection{DELKAPITTEL}
While this practice of businesses benefitting from or purchasing user data may be profitable, it can in some cases be a violation of user privacy. The handling and management of both public and private information has therefore been the subject of many studies in this era of Big Data (Buchanan, et al., 2006).  … (skriv mer om dette) 

\section{Conclusion}

% Kort oppsummering av hva du har kommet frem til. Oppsummer hvordan oppgaven svarte på motivation/spørsmål fra introduksjonen av oppgaven.


%\bibliography{refs}
%\bibliographystyle{plain}

\section{References}
Devore, J. L., \& Berk, K. L. (2018) Modern Mathematical Statistics with Application. New York Dordrecht Heidelberg London: Springer.\\ \\

Pedamkar, Priya. "Machine Learning vs Statistics." Asia's Largest Online Learning Platform, https://educba.com/machine-learning-vs-statistics/. Accessed Nov. 1st 2022. 
\\ \\

Kumar, Sanjeev. "AI Outperforms Humans in Question Answering." Medium - Stay curious., Mar 28 2018, https://medium.com/the-new-nlp/ai-outperforms-humans-in-question-answering-70554f51136b. Accessed Nov. 1st 2022.\\ \\

Lipeles, Aaron. "The Unreasonable Effectiveness of Data." towards data science, Oct 22, 2018, https://towardsdatascience.com/ai-ml-practicalities-the-unreasonable-effectiveness-of-data-c0bfd44c5057. Accessed Nov. 1st 2022.\\ \\

Yinhan Liu, Myle Ott, Naman Goyal, Jingfei Du, Mandar Joshi, Danqi Chen, Omer Levy, Mike Lewis, Luke Zettlemoyer, Veselin Stoyanov. "RoBERTa: A Robustly Optimized BERT Pretraining Approach." Cornell University - arXiv, https://arxiv.org/pdf/1907.11692.pdf.\\ \\

Alava, S., Frau-Meigs, D. \& Hassan, G., 2017. Youth and Violent Extremism on Social Media: Mapping the Research. s.l.:United NAtions Educational, Scientific and Cultural Organization.\\ \\

Arya, V., Sethi, D. \& Paul, J., 2019. Does digital footprint act as a digital asset? – Enhancing brand experience through remarketing. s.l.:International Journal of Information Management. \\ 

Buchanan, T., Paine, C., Joinson, A. N. \& Reips, U., 2006. Development of measures of online privacy concern and protection for use on the Internet. s.l.:Wiley Online Library. \\

Cate, F. H., 2008. Government data mining: The Need for a Legal Framework. s.l.:Harvard Civil Rights-Civil Liberties Law Review. \\

Etternavn1, F. \& Etternavn2, F., 2022. Et publisert tidskrift. Tidskrift 1, Oktober, pp. 1-12. \\

Jøsang, A., 2021. Informasjonssikkerhet - teori og praksis. s.l.:Universitetsforlaget. \\

Karas, S., 2002. Enhancing the privacy discourse: consumer information gathering as surveillance.. s.l.: J. Tech. L. \& Pol'y.. \\

Kerr, O., 2021. Buying Data and the Fourth Amendment. s.l.:Lawfare Institute. \\

Kaashoek, J. \& Santillana, M., 2020. COVID-19 positive cases, evidence on the time evolution of the epidemic or an indicator of local testing capabilities? A case study in the United States.. s.l.:SSRN. \\

McCausland, P. \& Schnecter, A., 2018. Trump-Linked Consultants Harvested Data from Millions on Facebook.. s.l.:NBCUniversal News Group. \\

Mulvenna, M. D., Anand, S. S. \& Büchner, A. G., 2000. Personalization on the Net using Web mining. s.l.:Communications of the ACM. \\

Munzert, S. et al., 2021. Tracking and promoting the usage of a COVID-19 contact tracing app. s.l.:Nature Human Behaviour. \\

Palos-Sanchez, P., Saura, J. R. \& Martin-Velicia, F., 2019. A study of the effects of programmatic advertising on users' concerns about privacy overtime. s.l.:Elsevier. \\

Personal Finance, 2019. What is a ChoicPoint Rrport?. [Internett]  
Available at: https://personal-finance.extension.org/what-is-a-choicepoint-report/ 
[Funnen 25 October 2022]. \\

Rathi, R., 2019. Effect of Cambridge Analytica’s Facebook ads on the 2016 US Presidential Election. s.l.:Towards Data Science. \\

Reporters Committee for Freedom of the Press, 2009. Out of sight, out of bounds: Article from the Spring 2009 issue of The News Media \& The Law. [Internett]  
Available at: https://www.rcfp.org/journals/the-news-media-and-the-law-spring-2009/out-sight-out-bounds/ 
[Funnen 25 October 2022]. \\

Russel, M. A., 2013. Mining the Social Web. 2. red. s.l.:O'Reilly Media. \\

Sinha, G. A., 2013. NSA surveillance since 9/11 and the human right to privacy. s.l.:s.n.\\ 

Sumpter, D. J. T., 2018. David J. T. Outnumbered: from Facebook and Google to Fake News and Filter-Bubbles — the Algorithms That Control Our Lives.. s.l.:Bloomsbury Sigma. \\

Testa-Avila, E., 2008. Practical Obscurity in the Digital Age: Public Records in the Private Sector. s.l.:UC Berkeley. \\

Trang, P. T., 2017. Personalized ads on Facebook: An effective marketing tool for online marketers. s.l.:Elsevier. \\

Vargo, C. J. \& Hopp, T., 2020. Fear, Anger, and Political Advertisement Engagement: A Computational Case Study of Russian-Linked Facebook and Instagram Content. s.l.:Journalism \& Mass Communication Quarterly. \\

Zuboff, S., 2015. Big other: surveillance capitalism and the prospects of an information civilization. s.l.:Journal of Information Technology. \\

Zuboff, S., 2019. The Age of Surveillance Capitalism: The Fight for a Human Future at the New Frontier of Power. s.l.:Profile Books. \\


\end{document}
